\documentclass[11pt,a4paper]{article}
\usepackage[margin=2.5cm]{geometry}
\usepackage{amsmath,amssymb}
\usepackage{physics}
\usepackage{booktabs}
\usepackage{hyperref}
\usepackage{parskip}

\newcommand{\kperp}{k_\perp}
\newcommand{\kpar}{k_\parallel}
\newcommand{\Tsys}{T_\mathrm{sys}}
\newcommand{\Pnoise}{P_\mathrm{N}}
\newcommand{\Psignal}{P_\mathrm{signal}}
\newcommand{\fnoise}{f_\mathrm{noise}}
\newcommand{\ffore}{f_\mathrm{fore}}
\newcommand{\nbl}{n_\mathrm{bl}}

\title{Theory: Thermal Noise in 21\,cm Interferometry\\[4pt]
\large From baselines to the 2D noise filter}
\author{Meeting notes}
\date{\today}

\begin{document}
\maketitle
\tableofcontents
\newpage

%======================================================================
\section{What does an interferometer measure?}
%======================================================================

Each pair of antennas (a \emph{baseline} with vector $\vb{b}$) measures one complex visibility:
\begin{equation}
    V(u,v,\nu) = \int I(l,m,\nu)\, B(l,m,\nu)\, e^{-2\pi i(ul + vm)}\, \dd{l}\,\dd{m}
\end{equation}
where $I$ is the sky brightness, $B$ is the primary beam, and $(l,m)$ are direction cosines on the sky.

The baseline coordinates in units of wavelengths are:
\begin{equation}
    u = \frac{b_\mathrm{east}}{\lambda}, \qquad
    v = \frac{b_\mathrm{north}}{\lambda}, \qquad
    \lambda = \frac{c}{\nu}
\end{equation}

A baseline of length $d$ probes angular scale $\theta \sim \lambda/d$. The visibility $V(u,v)$ is one Fourier coefficient of the sky brightness.

%======================================================================
\section{From baseline to wavenumber}
%======================================================================

The observed 21\,cm frequency maps to redshift:
\begin{equation}
    \nu_\mathrm{obs} = \frac{\nu_{21}}{1+z}, \qquad \nu_{21} = 1420.405\;\mathrm{MHz}
\end{equation}

Two independent Fourier axes connect observables to cosmological wavenumbers:

\paragraph{Transverse (angular $\to$ $\kperp$):}
\begin{equation}
    \kperp = \frac{2\pi\, u}{D_M(z)}
\end{equation}
where $D_M(z)$ is the comoving transverse distance (the paper calls this $X \equiv D_M$).

\paragraph{Radial (frequency $\to$ $\kpar$):}
\begin{equation}
    \kpar = \frac{2\pi\, \eta}{Y}, \qquad
    Y \equiv \frac{c\,(1+z)^2}{H(z)\,\nu_{21}} \quad [\mathrm{Mpc/MHz}]
\end{equation}
where $\eta$ is the Fourier dual of frequency (delay, in seconds).

Key point: $\kperp$ is set by \emph{which baseline} you use; $\kpar$ is set by \emph{which frequency channel} you observe.

%======================================================================
\section{Thermal noise per visibility}
%======================================================================

Each visibility measurement has independent Gaussian thermal noise with RMS:
\begin{equation}
    \sigma_\mathrm{vis} = \frac{\Tsys}{\sqrt{2\,\Delta\nu\,t_\mathrm{int}}}
\end{equation}
where:
\begin{itemize}
    \item $\Tsys$: system temperature [K], dominated by galactic synchrotron at EoR frequencies
    \item $\Delta\nu$: channel bandwidth [Hz]
    \item $t_\mathrm{int}$: integration time per visibility [s]
    \item Factor 2: two polarisations
\end{itemize}

The system temperature at EoR frequencies is modelled as (Eq.~19 of the paper):
\begin{equation}
    \Tsys(\nu) = 237 + 1.6\left(\frac{\nu}{300\;\mathrm{MHz}}\right)^{-5.23} \quad [\mathrm{K}]
\end{equation}
At $z=8$ ($\nu = 158$\,MHz): $\Tsys \approx 700$\,K.

This noise is \textbf{independent} between different baselines and different frequency channels. This means the noise is \textbf{white} in $(u,v,\nu)$ space --- a crucial property.

%======================================================================
\section{From visibility noise to power spectrum noise}
%======================================================================

The thermal noise power spectrum $\Pnoise$ measures how much noise power there is per Fourier mode. For HERA, this is given by Eq.~14 of the paper (Parsons+ 2014 formalism):
\begin{equation}\label{eq:pnoise}
    \boxed{
    \Pnoise(\kperp) = \frac{\Tsys^2\;\Omega_p^2/\Omega_{pp}\; X^2\, Y}{t_\mathrm{int}\; N_\mathrm{pol}\; \nbl(u)}
    }
\end{equation}

\begin{table}[h]
\centering
\begin{tabular}{llr}
\toprule
Factor & Physical meaning & Value at $z=8$ \\
\midrule
$\Tsys^2$ & Noise amplitude squared & $(700\;\mathrm{K})^2$ \\
$\Omega_p^2/\Omega_{pp}$ & Beam correction ($\approx 2\,\Omega_p$ for Gaussian) & $\sim 0.037$\,sr \\
$X^2 = D_M^2$ & Angles $\to$ comoving Mpc$^2$ & $(6500\;\mathrm{Mpc})^2$ \\
$Y$ & Frequency $\to$ comoving Mpc/MHz & $\sim 12$\,Mpc/MHz \\
$t_\mathrm{int}$ & Integration time (more $\to$ less noise) & $200\;\mathrm{h}$ \\
$N_\mathrm{pol}$ & Polarisations & 2 \\
$\nbl(u)$ & Baseline density at this $u$ & varies with $\kperp$ \\
\bottomrule
\end{tabular}
\caption{Factors in the HERA noise power spectrum.}
\end{table}

For SKA1-Low, the noise takes a slightly different form (Eq.~17, McQuinn+ 2006):
\begin{equation}
    \Pnoise^\mathrm{SKA}(\kperp) = \frac{\lambda^2\;\Tsys^2\; X^2\, Y}{t_\mathrm{int}\; \nbl(u)\; A_e}
\end{equation}
where $A_e$ is the effective collecting area per station.

%======================================================================
\section{Why $\Pnoise$ depends on $\kperp$ but NOT on $\kpar$}
%======================================================================

This is the single most important point for understanding the 2D filter:

\begin{enumerate}
    \item $\kperp$ is set by \emph{which baseline} you use $\to$ $\nbl(u)$ varies with baseline length $\to$ $\Pnoise$ varies with $\kperp$
    \item $\kpar$ is set by \emph{frequency} $\to$ thermal noise is \textbf{white in frequency} (independent between channels) $\to$ after FFT to $\kpar$ space, the noise power is \textbf{flat}
\end{enumerate}

\begin{equation}
    \Pnoise = \Pnoise(\kperp) \quad \text{--- independent of } \kpar
\end{equation}

This is why the 2D filter has \textbf{vertical band structure}: at a given $\kperp$, the filter value is the same for all $\kpar$ (outside the wedge).

%======================================================================
\section{The baseline density $\nbl(u)$}
%======================================================================

$\nbl(u)$ is the number of baseline pairs sampling a given $(u,v)$ annulus. It depends entirely on the antenna layout:

\paragraph{HERA} (compact hexagonal, 350 dishes, 14\,m spacing):
\begin{itemize}
    \item Many short baselines (14, 28, 42\,m\ldots) $\to$ dense coverage at small $u$ (low $\kperp$)
    \item Few long baselines $\to$ sparse at large $u$ (high $\kperp$)
    \item $\nbl$ peaks at $u \sim 50\,\lambda$ ($\kperp \sim 0.05\;h\,\mathrm{Mpc}^{-1}$)
\end{itemize}

\paragraph{SKA-Low} (extended, 512 stations, up to 65\,km):
\begin{itemize}
    \item Few short baselines $\to$ sparse at small $u$ (low $\kperp$)
    \item Many intermediate baselines $\to$ peak at $u \sim 300\,\lambda$
    \item Result: noise suppresses \textbf{small-scale} (high $\kperp$) modes
\end{itemize}

This explains the key visual difference between HERA and SKA in the paper's Fig.~3: HERA's noise kills high $\kperp$ (right side dark), while SKA's noise kills low $\kperp$ (left side dark).

%======================================================================
\section{The Wiener noise filter}
%======================================================================

The optimal linear filter for extracting signal from noise is the Wiener filter (Eq.~13):
\begin{equation}\label{eq:wiener}
    \boxed{
    \fnoise(\kperp) = \frac{\Psignal(k)}{\Psignal(k) + \Pnoise(\kperp)}
    }
\end{equation}

Limiting cases:
\begin{align}
    \Psignal \gg \Pnoise &\implies \fnoise \to 1 \quad \text{(mode survives)} \\
    \Psignal \ll \Pnoise &\implies \fnoise \to 0 \quad \text{(mode lost)} \\
    \Psignal = \Pnoise &\implies \fnoise = 0.5 \quad \text{(crossover)}
\end{align}

Since $\Pnoise$ depends only on $\kperp$, the filter has the same value for all $\kpar$ at a given $\kperp$.

Note: $\Psignal(k)$ can depend on $k = \sqrt{\kperp^2 + \kpar^2}$, so the filter is not \emph{perfectly} flat in $\kpar$ --- but the dominant variation is through $\kperp$ via $\nbl(u)$.

%======================================================================
\section{The foreground wedge}
%======================================================================

Bright astrophysical foregrounds (galactic synchrotron, point sources) are spectrally smooth and have power only at low $\kpar$. However, the interferometer's chromatic response (beam width $\propto \lambda/D$ changes with frequency) leaks foreground power into a wedge-shaped region in $(\kperp, \kpar)$ space.

The wedge boundary (Eq.~12):
\begin{equation}\label{eq:wedge}
    \boxed{
    \ffore(\kperp, \kpar) = \Theta\!\left(|\kpar| - m(z)\,\kperp\right)
    }
\end{equation}
where $\Theta$ is the Heaviside step function and the slope is:
\begin{equation}
    m(z) = \frac{H(z)\,D_M(z)}{c\,(1+z)}
\end{equation}

At $z=8$: $m \approx 0.32$ (cosmological value). In practice, a more conservative slope $m = 1.0$ (the ``horizon limit'') is often used. The \texttt{tuesday} package uses $m=1.0$ by default.

Modes \emph{inside} the wedge ($|\kpar| < m\,\kperp$) are contaminated and zeroed. The safe region outside is called the \textbf{EoR window}.

%======================================================================
\section{The combined 21\,cm filter}
%======================================================================

The total filter applied to the 21\,cm field (Eq.~18):
\begin{equation}\label{eq:total}
    \boxed{
    f^{21\mathrm{cm}}(\kperp, \kpar) = \ffore(\kperp, \kpar) \times \fnoise(\kperp)
    }
\end{equation}

Reading the 2D plot:
\begin{itemize}
    \item \textbf{Bottom-left triangle}: killed by wedge ($\ffore = 0$)
    \item \textbf{Left column} (low $\kperp$, outside wedge): killed by noise for HERA (few short baselines at very small $u$), surviving for SKA
    \item \textbf{Right column} (high $\kperp$, outside wedge): killed by noise for SKA (sparse long baselines), surviving for HERA's core
    \item \textbf{Middle band}: the \textbf{EoR window} --- where signal survives both noise and foregrounds
\end{itemize}

%======================================================================
\section{Computing $\Pnoise$: three approaches}
%======================================================================

\subsection{A. Paper-analytical (Eq.~\ref{eq:pnoise})}

Use the formula directly with an analytic approximation for $\nbl(u)$. The implementation in \texttt{fig3\_filter.py} uses a log-normal fit tuned to match the paper's Fig.~3.

\textbf{Pros:} no simulation needed, fast, perfectly smooth, reproducible.\\
\textbf{Cons:} $\nbl(u)$ is approximate (not the real antenna positions).

\subsection{B. Tuesday-analytical}

Use \texttt{compute\_thermal\_rms\_uvgrid()} from the \texttt{tuesday} package, which computes $\sigma_\mathrm{noise}$ per UV cell from the actual HERA antenna positions and observation parameters. Then $\Pnoise = \sigma^2$.

\textbf{Pros:} uses real antenna layout, accounts for rotation synthesis.\\
\textbf{Cons:} only non-zero where baselines exist (sparse UV grid), requires careful binning.

\subsection{C. Noise realizations}

Pass a zero-signal box through \texttt{observe\_coeval} $\to$ pure noise output. Compute its power spectrum. Optionally average over $N$ realizations.

\textbf{Pros:} includes all effects (beam, gridding, interpolation).\\
\textbf{Cons:} single realization has variance $\sim 1/N_\mathrm{modes}$ per bin.

The variance of the power spectrum estimate in a bin:
\begin{equation}
    \frac{\sigma(P)}{P} \sim \frac{1}{\sqrt{N_\mathrm{modes}}}
\end{equation}
\begin{itemize}
    \item At \textbf{high $\kperp$}: many modes (many integer pairs $(n_x, n_y)$ give similar $\sqrt{n_x^2 + n_y^2}$) $\to$ single realization is fine
    \item At \textbf{low $\kperp$}: few modes $\to$ large scatter $\to$ dark patches in the filter
\end{itemize}
Averaging over $N$ realizations reduces scatter by $\sqrt{N}$, converging to the analytical result.

%======================================================================
\section{Order of operations for the kSZ$^2\times$21cm$^2$ pipeline}
%======================================================================

The paper (Section~3) emphasises the following order:
\begin{enumerate}
    \item Extract a redshift chunk from the 21\,cm light cone
    \item Fourier transform the 3D chunk
    \item Apply filters: $f^{21\mathrm{cm}}$ for 21\,cm, $f_\mathrm{kSZ}$ for kSZ
    \item Inverse FFT back to real space
    \item \textbf{Square} the fields in configuration space: $(\delta T_{21})^2$, $(\delta T_\mathrm{kSZ})^2$
    \item \textbf{Then} project 21\,cm to 2D by integrating along LOS
    \item Compute cross-power spectrum between squared fields
\end{enumerate}

If you project \emph{before} squaring, the radial modes that survived filtering are lost in the projection. Squaring first preserves the 3D mode information.

This ordering is specific to the kSZ$^2\times$21cm$^2$ cross-correlation. For a direct (unsquared) cross-correlation, this issue does not apply.

%======================================================================
\section{Summary of key numbers at $z=8$}
%======================================================================

\begin{table}[h]
\centering
\begin{tabular}{lr}
\toprule
Quantity & Value \\
\midrule
Observed frequency $\nu_\mathrm{obs}$ & 158\,MHz \\
Wavelength $\lambda$ & 1.9\,m \\
System temperature $\Tsys$ & $\sim 700$\,K \\
Comoving distance $D_M = X$ & $\sim 6500$\,Mpc \\
LOS conversion $Y$ & $\sim 12$\,Mpc/MHz \\
Wedge slope $m(z)$ & 0.32 (cosmological), 1.0 (horizon limit) \\
HERA dishes & 350, diameter 14\,m \\
Integration time & 200\,h (paper) / 1080\,h (our setup) \\
\bottomrule
\end{tabular}
\end{table}

\end{document}
